\documentclass[11pt,a4paper]{article}
\usepackage[margin=24mm]{geometry}
\usepackage{amsmath,amssymb,amsthm}
\setlength{\emergencystretch}{2em}
\newtheorem{lemma}{Lemma}
\newtheorem{theorem}[lemma]{Theorem}
\newcommand{\gen}[1]{\langle#1\rangle}
\title{Prefix orbits from a singleton end}
\author{Human Compression project: L044 continuation}
\date{}
\begin{document}\maketitle
Let $G=\gen{\rho_0,\ldots,\rho_{r-1}}\le\operatorname{Alt}(\Omega)$
be transitive with a connected string C-marking, and suppose every
maximal omission is intransitive. Suppose $G_0$ has the singleton
orbit $\{a\}$. Put $P_k=G_{\le k}$ and $C_k=a^{P_k}$.
The other $G_0$-orbit is understood to be its entire complement,
as in the complete singleton-end interface.

\begin{lemma}[Strict growth without an alternating kernel]\label{growth}
One has $|C_0|=2$ and $C_{k-1}\subsetneq C_k$ for
$1\le k<r$. Every mark of index at least $k+2$ fixes
$C_k$ pointwise. In particular $|C_k|\ge k+2$.
\end{lemma}
\begin{proof}
Every mark except $\rho_0$ fixes $a$. Transitivity makes
$a^{\rho_0}\ne a$, proving the first assertion. A mark of
index at least $k+2$ centralizes $P_k$ and fixes $a$,
so fixes its entire $P_k$-orbit. If $C_k=C_{k-1}$,
the marks up to index $k$ preserve this set and all later
marks fix it pointwise by the same observation applied to
$P_{k-1}$. It would be $G$-invariant and hence equal to
$\Omega$. But $P_{k-1}\le G_k$, so this would make that
omission transitive, a contradiction. The size bound follows.
No claim of a pure alternating kernel is used or implied.
\end{proof}

\begin{lemma}[The actual sides of a later cut]\label{sides}
If a colour $j\ge2$ has exactly one ambient-eligible edge,
its two omission orbits are $C=C_{j-1}$ and
$A=\Omega\setminus C$. The suffix $G_{>j}$ fixes $C$
pointwise. Every shared label at this cut is below $j$.
\end{lemma}
\begin{proof}
The suffix centralizes $P_{j-1}$ and fixes $a$, so fixes
$C$. Thus $C$ is exactly the $G_j$-orbit of $a$.
A single crossing edge can connect all omission orbits only
if there are exactly two. The suffix cannot act on both,
giving the final assertion.
\end{proof}

For the next theorem retain the complete acquired support,
small-tail and critical-dihedral proofs, including their explicit
smaller-degree alternating induction. Assume $n=|\Omega|\ge12$
and the contrary range $r>\lfloor(n-1)/2\rfloor$.
These are precisely the hypotheses under which the complete
critical singleton reduction gives $n=2r\ge14$ and
$G_0=\operatorname{Alt}(\Omega\setminus\{a\})$.

\begin{theorem}[The first shared indices at a singleton end]
Suppose a later unique cut $j\ge2$ has a first shared label $g$.
Then $g\ge3$ and
\[
             |\operatorname{supp}(P_{g-1})|\le2g+1.       \tag{1}
\]
The small-index exclusions use the actual singleton point
and the complete critical-dihedral-tail proof, without assuming
that the cut adjacent to a supported head is itself unique.
\end{theorem}
\begin{proof}
Reverse the entire string. Its crossing is $i=r-1-j$;
the last shared label is $h=r-1-g$. If $i>0$, the group
preceding this crossing is $G_{>j}$ and fixes $C$ pointwise
by Lemma~\ref{sides}. Both omission orbits are nontrivial:
$C$ has at least three points, and a shared label acts on $A$.
Thus the complete interior support alternative applies,
using its transitive or intransitive case according to the
actual action of $G_{>j}$ on $A$. It excludes $h=r-1$
in the contrary range. If $i=0$, the complete endpoint
shared-last-label theorem gives the same exclusion.
Consequently $g>0$.

\newpage
All labels below $g$ are unshared. They act on the same
omission orbit, since a change of sides among consecutive
nonidentity marks would disconnect the noncommutation diagram.
This side is $C$, because $\rho_0$ moves $a$.
The entire reversed tail is therefore exactly the actual
supported prefix $P_{g-1}$. The complete support bound,
interior or endpoint as appropriate, gives
\[
 n-|\operatorname{supp}(P_{g-1})|\ge2h+1.
\]
Substituting $n=2r$ and $h=r-1-g$ proves (1).
For $g=1$ this bounds the support of the nonidentity even
involution $\rho_0$ by three, impossible. Hence $g\ge2$.

Suppose $g=2$. The support of $P_1$ has at most five points,
and its orbit $C_1$ has at least three by Lemma~\ref{growth}.
If $P_1$ is intransitive on its support, there is exactly one
other nontrivial orbit $E$, with $|C_1|=3$ and $|E|=2$.
Write $C_1=\{a,b,c\}$, where $b=a^{\rho_0}$.
The two restrictions to $C_1$ are $(a\ b)$ and $(b\ c)$:
$\rho_1$ fixes $a$, and the displayed orbit is transitive.
Global evenness then forces both marks to exchange the pair
$E=\{u,v\}$. Both fix every other point.

The mark $\rho_2$ fixes $a$ and centralizes
$\rho_0=(a\ b)(u\ v)$, so fixes $b$ and preserves $E$.
Every mark of index at least three centralizes $P_1$.
Its two nontrivial orbits have distinct sizes three and two,
so such a mark preserves $E$ as well. Hence all of $G$
preserves this proper pair, contrary to transitivity.
Thus $P_1$ is transitive on its support.

A connected even rank-two marking cannot act transitively
on at most four points: degrees below four have no nonidentity
even involutions, and the involutions of $A_4$ commute.
Its support therefore has exactly five points. The two
generators are double transpositions. Their product has
order two, three or five by its cycle type in $A_5$;
orders two and three give groups of orders at most four
and six, neither transitive on five points. Its order is
therefore five, and the actual group is the natural $D_{10}$.
Its two distinct reflection generators have the standard ordered
marking after a change of vertex names: in the affine action on
$\mathbb F_5$, a translation and nonzero scaling normalize
the two reflection parameters.
Apply now the complete critical-dihedral-tail lemma at the
original cut, in the already established reversed orientation.
Its supported tail past $h=r-3$ is exactly this actual $P_1$,
of degree five, and its fixed deficit is
$n-5=2r-5=2h+1$. If the reversed crossing is interior,
its preceding subgroup fixes the other orbit $C$ pointwise
by Lemma~\ref{sides}; the endpoint version has no such
preceding subgroup. These are precisely the two cases of the
retained lemma, which excludes both for $n\ge12$.
Thus $g=2$ is impossible as well. No claim that the
separate original colour $2$ has a unique eligible edge, or
that its opposite suffix is transitive, has been inserted.
\end{proof}

\paragraph{Attribution and boundary.}
This develops the remaining singleton route in the source
Propositions~5.28--5.32. The complete reversed support and
critical-dihedral proofs are retained with their attribution
to Cameron--Fernandes--Leemans--Mixer. The strict prefix growth
uses the distinguished point fixed by every later mark; its
argument resembles the acquired head-prefix growth but does
not import that theorem's alternating kernel.
The exterior pair exclusion above is a literal permutation
argument. First shared indices $g\ge3$, unshared later cuts,
and the complete singleton and principal source assembly remain.
No whole-original, question, numerical or formal credit is claimed.
\end{document}
